% ======================================================================
%  Kapitel 3 - Datenaugmentierung und synthetische Datengenerierung
%
%  Zur Einbindung in eine bestehende Arbeit nur den Kapitelinhalt
%  uebernehmen; die Makros (\Dom, \Hcls, \HdH, \IoU, \px, \absatz,
%  \mAPfifty) stammen aus der Praeambel der Arbeit.
%
%  Die Gleichungsnummern (3.1)-(3.13) sind gegenueber der vorherigen
%  Fassung unveraendert. Neu hinzugekommen ist (3.14) in Abschnitt 3.4;
%  die Definition der HdH-Divergenz steht als unnummerierte Gleichung,
%  damit sich keine Nummer verschiebt.
%
%  Zitiert wird mit biblatex (\parencite, \textcite); alle Quellen stehen
%  in literatur.bib.
% ======================================================================

%\documentclass[11pt,a4paper]{report}



%\begin{document}

\chapter{Datenaugmentierung und synthetische Datengenerierung}
\label{kap:augtheorie}

Datenaugmentierung bezeichnet die Erzeugung zusätzlicher Trainingsbeispiele aus
vorhandenen Daten durch Transformationen, die die für die Lernaufgabe relevante
Beschriftung erhalten \parencite{shorten2019}. Ziel ist es, die
Verteilung der im Training beobachteten Beispiele zu erweitern und das Modell
gegenüber Variationen robust zu machen, die bei der späteren Anwendung
auftreten können.

Im vorliegenden Anwendungsfall wird dieser Begriff weiter gefasst. Die
Trainingsdaten werden nicht lediglich durch Transformationen real aufgenommener
Bilder vergrößert. Stattdessen werden die Trainingsbilder überwiegend
synthetisch aus wenigen realen Ausgangsobjekten und Hintergründen erzeugt --
ein Vorgehen, das für die Detektion bereits in Innenraumszenen erprobt wurde
\parencite{georgakis2017}. Die
zentrale Frage ist daher nicht nur, welche Transformationen auf vorhandene
Bilder angewendet werden sollen, sondern welche Verteilung synthetischer Szenen
erzeugt werden muss, damit sie die für die Inferenz relevante Domäne
hinreichend abbildet.

Für die Notation gilt in diesem Kapitel durchgehend: lateinische Großbuchstaben
wie $D$ bezeichnen endliche Datensätze, kalligraphische Symbole wie
$\Dom{S}^{X}$ bezeichnen Wahrscheinlichkeitsverteilungen.

% ----------------------------------------------------------------------
\section{Begriff, Formalisierung und Wirkmechanismen}
\label{sec:aug-formalisierung}

Sei $D=\{(x_i,y_i)\}_{i=1}^{n}$ ein Trainingsdatensatz mit Eingabebildern $x_i$
und Beschriftungen $y_i$, bei der Objektdetektion eine Menge von Boxen mit
Klasseninformation. Ein Modell $f$ bildet ein Bild auf eine Vorhersage ab, und
eine Verlustfunktion $L$ misst, wie weit eine Vorhersage von der Beschriftung
abweicht. Sei weiterhin $T$ eine Menge zulässiger Transformationen,
$p(t)$ eine Wahrscheinlichkeitsverteilung über diese Menge und $\tau_t$ die zu
$t \in T$ gehörende Transformation der Beschriftung. Das durch Augmentierung
entstehende empirische Risiko lautet dann
\begin{equation}
  \hat R_{\mathrm{aug}}(f)
  =
  \frac{1}{n}\sum_{i=1}^{n}
  \mathbb{E}_{t\sim p(t)}
  \Bigl[
    L\bigl(f(t(x_i)),\,\tau_t(y_i)\bigr)
  \Bigr].
  \label{eq:risk-aug}
\end{equation}

Der Unterschied zum gewöhnlichen empirischen Risiko liegt allein im
Erwartungswert über $t$: Bewertet wird nicht jedes Trainingsbeispiel in seiner
ursprünglichen Form, sondern eine Verteilung von Varianten -- die Sichtweise der
Vicinal Risk Minimization \parencite{chapelle2001}. Da $T$ im Allgemeinen
kontinuierlich ist, wird der Erwartungswert praktisch über endlich viele
Ziehungen aus $p(t)$ geschätzt; die Zahl physisch gespeicherter Bilder ist
deshalb nicht identisch mit der Zahl der im Training gesehenen Varianten.

Von den vier Größen in Gleichung~\eqref{eq:risk-aug} -- Modell $f$,
Verlust $L$, Transformationsverteilung $p(t)$ und Labeltransformation
$\tau_t$ -- ist für diese Arbeit vor allem $p(t)$ maßgeblich. Sie bestimmt, wie häufig und wie stark die einzelnen
Transformationen wirken, und ist damit keine Implementierungsentscheidung,
sondern eine \emph{Modellierung der erwarteten Variation der Zieldomäne}. Bei
überwiegend synthetischen Trainingsdaten verschiebt sich ihre Bedeutung noch
weiter: Dann beschreibt $p(t)$ nicht mehr nur die Variation um vorhandene
Beispiele herum, sondern einen wesentlichen Teil der beobachteten
Datenverteilung selbst. Die Labeltransformation $\tau_t$ ist dabei die stille
Fehlerquelle -- eine Helligkeitsänderung lässt die Box unverändert, eine
Rotation dagegen verlangt, dass die achsparallele Box aus der transformierten
Objektgeometrie neu bestimmt wird. Ist der Einfluss auf die Beschriftung nicht
exakt bekannt, entsteht zusätzliches Labelrauschen, und bei kleinen Objekten
betrifft bereits eine geringe geometrische Abweichung einen vergleichsweise
großen Anteil der Objektfläche.

\absatz{Invarianz, nicht Vielfalt.} Augmentierung zeigt dem Modell Variationen,
unter denen die Zielentscheidung invariant bleiben soll, wirkt über zufällige
photometrische Störungen als Regularisierung gegen die Festlegung auf einzelne
exakt reproduzierbare Bildmerkmale \parencite{bishop1995} und erweitert den
belegten Bereich des Eingaberaums. Die zusätzlichen Beispiele sind jedoch keine
unabhängigen Beobachtungen der Objektklasse: Ihre Information stammt weiterhin
aus den ursprünglichen Daten und den Annahmen des Generators. Daraus folgt das
Prinzip, das die gesamte Pipeline trägt:

\begin{quote}
  Eine Transformation sollte nur in dem Maß als labelerhaltend behandelt werden,
  in dem die entsprechende Variation auch in der Zieldomäne auftreten kann.
\end{quote}

Wird sie darüber hinaus angewendet, entfernt sie diskriminative Information,
statt Robustheit zu erzeugen. Für die Rohrstirnflächen ist beispielsweise eine
große Variation der Orientierung plausibel, eine vollständige
Rotationsinvarianz für andere Detektionsaufgaben dagegen nicht.

\section{Verfahrensüberblick}
\label{sec:aug-verfahren}

Datenaugmentierung umfasst eine große Anzahl unterschiedlicher Verfahren. Für
die vorliegende Arbeit sind insbesondere geometrische Transformationen,
photometrische Transformationen, Verdeckungsstrategien und kompositionelle
Verfahren relevant; Tabelle~\ref{tab:taxonomie} im Anhang ordnet die
Verfahrensklassen ihrem Einsatz im vorliegenden Fall zu.

Geometrische Transformationen verändern die Lage oder Form von Objekten. Dazu
zählen insbesondere Translation, Skalierung, Spiegelung, Rotation, Scherung und
Perspektivtransformation. Photometrische Transformationen verändern dagegen die
Bildintensitäten, ohne die Objektgeometrie grundsätzlich zu verändern. Beispiele
sind Helligkeits- und Kontraständerungen, Farbverschiebungen, Rauschen,
Unschärfe und Kompression.

Eine weitere Klasse bilden Verfahren zur künstlichen Verdeckung von
Bildbereichen. Cutout \parencite{devries2017cutout}, Random Erasing
\parencite{zhong2020erasing} und Hide-and-Seek \parencite{singh2017hideseek}
erzeugen beispielsweise teilweise verdeckte Eingaben, um das Modell gegenüber fehlender Bildinformation
robuster zu machen. Für die vorliegende Anwendung ist jedoch entscheidend, dass
Verdeckung kein rein künstliches Störphänomen ist. Die Zielobjekte werden in der
realen Szene tatsächlich teilweise durch andere Profile, Gitterstrukturen und
Schatten verdeckt.

Daher wird Verdeckung nicht primär durch zufälliges Löschen rechteckiger
Bildbereiche modelliert. Stattdessen werden physikalisch plausible Objekte als
Okkluder in die Szene eingefügt. Dadurch können gleichzeitig die Geometrie, die
sichtbaren Objektbereiche und die Interaktion zwischen Vorder- und Hintergrund
kontrolliert werden.

\absatz{Augmentierung und Synthese.} Bei klassischer Datenaugmentierung bildet ein reales Trainingsbild den
Ausgangspunkt. Transformationen erzeugen daraus weitere Varianten:
\[
  x_i \;\rightarrow\; t(x_i).
\]
Die Information über die reale Bilddomäne stammt somit weiterhin aus dem
ursprünglichen Bild.

Bei der in dieser Arbeit verwendeten synthetischen Datengenerierung wird die
Szene dagegen aus mehreren Komponenten zusammengesetzt. Dadurch wird aus der
Augmentierung im engeren Sinne eine Form der prozeduralen Datengenerierung: Die
Vielfalt der Trainingsdaten entsteht nicht mehr aus einer großen Menge realer
Ausgangsbilder, sondern aus der Verteilung $p(t)$ selbst. Deren Struktur ist
damit Gegenstand der Modellierung und verdient eine genauere Betrachtung.

Ein Durchlauf zieht die Parameter in drei Stufen. Zuerst die szenenweiten
Größen $t^{\mathrm{S}}$ -- Hintergrund, Maßstab, Ausschnitt, globale
Bildeffekte --, daraus die Zahl $K$ der Objektcluster. Für jeden Cluster $k$
folgt ein eigener Vektor $c_k$ mit Basisgröße, Gitterabständen, eingeschlossenem
Winkel und Clusterzentrum, daraus die Zahl $m_k$ seiner Instanzen. Erst auf der
dritten Stufe erhält jede Instanz $j$ des Clusters $k$ ihren eigenen
Parametervektor $t^{(k,j)}$. Die Verteilung faktorisiert entsprechend als
\begin{equation}
  p(t)
  =
  p\bigl(t^{\mathrm{S}}\bigr)\,
  p\bigl(K \mid t^{\mathrm{S}}\bigr)
  \prod_{k=1}^{K}
  \Bigl[\,
    p\bigl(c_k \mid t^{\mathrm{S}}\bigr)\,
    p\bigl(m_k \mid c_k\bigr)
    \prod_{j=1}^{m_k} p\bigl(t^{(k,j)} \mid c_k\bigr)
  \,\Bigr].
  \label{eq:faktorisierung}
\end{equation}

Die Instanzparameter sind damit \emph{bedingt} unabhängig gegeben den Cluster,
nicht unabhängig. Diese Unterscheidung ist keine Formalie, sondern die formale
Fassung von Anforderung~A3: Objektgröße und Hintergrundmaßstab müssen gekoppelt
sein, weil Stirnflächen und Rohrschäfte derselben Aufnahme annähernd im selben
Maßstab abgebildet werden; und innerhalb eines Clusters streut jede Instanz nur
um eine gemeinsame Basisgröße, weil gestapelte Profile derselben Lage in der
Regel gleich groß sind. Über $t^{\mathrm{S}}$ und $c_k$ gemittelt sind die
Instanzgrößen daher korreliert -- ein Generator, der sie unabhängig zöge,
erzeugte Größenkombinationen, die in realen Aufnahmen kaum vorkommen. Wie groß
eine Stirnfläche im Verhältnis zu den Rohrschäften des Hintergrunds ist, bleibt
dabei eine Annahme: Gleichung~\eqref{eq:zoom} setzt beide Durchmesser
vereinfachend gleich, die seitlichen Rahmenrohre, deren Stirnflächen detektiert
werden, müssen aber nicht denselben Durchmesser haben wie die Rohre, deren
Schäfte im Hintergrund zu sehen sind. Vereinzelt platzierte Profile, die
keinem Cluster angehören, bilden den Sonderfall $m_k = 1$ und hängen unmittelbar
von $t^{\mathrm{S}}$ ab.

Da zudem $K$ und alle $m_k$ selbst gezogen werden, besitzt $t$ keine feste
Dimension: Der Träger von $p$ ist eine Vereinigung von Räumen unterschiedlicher
Dimension und nicht Teil eines Raums fester Dimension.

Jede einzelne optionale Transformation $\nu$ ist ihrerseits eine Mischung aus
Identität und Wirkung,
\begin{equation}
  p(t_\nu)
  =
  (1-\pi_\nu)\,\delta_{\mathrm{id}}
  \;+\;
  \pi_\nu\, q_\nu ,
  \label{eq:mischung}
\end{equation}
mit dem Anwendungsanteil $\pi_\nu$, der Parameterverteilung $q_\nu$ über den
zulässigen Bereich und der Verteilung $\delta_{\mathrm{id}}$, die mit Sicherheit
die Identität liefert, das Bild also unverändert lässt. Eine Transformation ist
folglich erst durch \emph{zwei}
Angaben bestimmt -- wie häufig sie wirkt und wie stark. Diese Doppelung
verlangt Anforderung~A5; Tabelle~\ref{tab:parameter} führt sie spaltenweise
auf.

\absatz{Nicht gewählte Wege.} Vier naheliegende Alternativen werden nicht verfolgt. Die ersten beiden
betreffen die Frage, \emph{wie} die Transformationsverteilung bestimmt wird, und
scheitern am selben Umstand: Beide setzen eine Menge realer Daten voraus, die
hier nicht existiert. Die beiden übrigen setzen an der Objektgröße an und werden
aus jeweils eigenen Gründen nicht benötigt.

Eine \emph{automatisierte Policy-Suche} wie AutoAugment \parencite{cubuk2019}
oder RandAugment \parencite{cubuk2020} optimiert Augmentierungsparameter gegen
eine Validierungsleistung -- hier hieße das, die Pipeline auf wenige
Einzelbilder anzupassen.

Eine \emph{generative Modellierung} der Objekte, etwa durch Diffusionsmodelle \parencite{trabucco2024},
benötigte eine hinreichend große und repräsentative Menge realer Beispiele, um
die Zielverteilung überhaupt zu erfassen. Bei neun Ausgangsobjekten ist das
nicht gegeben. Der explizit parametrisierte Generator bietet stattdessen die
hier benötigten Eigenschaften: unabhängige Kontrolle über Geometrie und Photometrie sowie
Labels, die direkt aus der bekannten Szenengeometrie folgen. Die Entscheidung
ist damit keine Aussage über die Leistungsfähigkeit generativer Modelle, sondern
folgt aus den Anforderungen an Kontrolle, Labelgenauigkeit und
Nachvollziehbarkeit.

Eine \emph{Kachelung des Eingabebildes}, wie sie Slicing Aided Hyper Inference
\parencite{akyon2022sahi} umsetzt, zerlegt das Bild in überlappende Ausschnitte
und wertet jeden einzeln aus. Kleine Objekte werden dadurch relativ zum
Netzeingang größer -- die Kachelung wirkt also auf dieselbe Größe
$d_{\mathrm{eff}}/S$ wie eine höhere Inferenzauflösung, nur über einen anderen
Weg. Dieser Weg ist für die vorliegende Aufgabe der ungünstigere: Die Objekte
stehen in dichten Clustern, sodass Schnittkanten regelmäßig durch
zusammenhängende Gruppen liefen, und das Zusammenführen der Teilergebnisse
verlangt eine Entscheidung über mehrfach detektierte Objekte an den
Kachelgrenzen -- also genau den Unterdrückungsschritt, auf den der verwendete
Ende-zu-Ende-Kopf verzichtet. Die Auflösung selbst zu erhöhen ist der
unmittelbarere Eingriff und wird in Kapitel~\ref{kap:auswertung} als solcher
vermessen.

Eine \emph{gezielte Überabtastung kleiner Objekte} \parencite{kisantal2019}
schließlich antwortet darauf, dass die Größenverteilung eines \emph{gegebenen}
Datensatzes nicht veränderbar ist. Diese Einschränkung besteht hier nicht: Die Größenverteilung ist ein Parameter des Generators und wird nach
Anforderung~A2 unmittelbar auf den Zielbereich kalibriert.

% ----------------------------------------------------------------------
\section{Besonderheiten bei der Objektdetektion}
\label{sec:aug-detektion}

Datenaugmentierung für Objektdetektion unterscheidet sich von der Augmentierung
bei Klassifikationsaufgaben insbesondere dadurch, dass eine Transformation nicht
nur das Bild, sondern auch die räumlichen Beschriftungen verändern muss.

Bei der Klassifikation kann für viele labelerhaltende Transformationen
\[
  \tau_t(y)=y
\]
angenommen werden. Bei der Detektion gilt dies im Allgemeinen nicht. Wird ein
Objekt beispielsweise um $(\Delta x,\Delta y)$ verschoben, müssen auch die
Boxkoordinaten entsprechend verschoben werden.

Für die vorliegende Aufgabe entstehen daraus zwei besondere Anforderungen.
Erstens müssen geometrische Transformationen möglichst exakt auf die
Beschriftungen übertragen werden. Zweitens muss die erzeugte Objektgröße mit der
räumlichen Auflösung des verwendeten Detektors vereinbar sein.

\absatz{Achsparallele Boxen und Rotation.}\label{subsec:aabb-rotation} Die Zielobjekte werden mit achsparallelen Bounding Boxes beschrieben. Diese
Darstellung ist für die verwendete Detektionsaufgabe zweckmäßig, führt jedoch bei
geometrischen Transformationen zu einer zusätzlichen Approximation.

Betrachtet werde zunächst ein Quadrat mit Seitenlänge $d$, das um einen Winkel
$\theta$ rotiert wird. Die Seitenlänge der kleinsten achsparallelen Bounding Box
des rotierten Quadrats beträgt
\begin{equation}
  d_{\mathrm{box}}
  =
  d\bigl(|\cos\theta|+|\sin\theta|\bigr).
  \label{eq:dbox}
\end{equation}

Wie gut zwei Boxen $B_1$ und $B_2$ übereinstimmen, misst die \emph{Intersection
over Union}, der Quotient aus Schnitt- und Vereinigungsfläche
\parencite{everingham2010}:
\begin{equation}
  \IoU(B_1, B_2)
  =
  \frac{|B_1 \cap B_2|}{|B_1 \cup B_2|}
  =
  \frac{|B_1 \cap B_2|}{|B_1| + |B_2| - |B_1 \cap B_2|} ,
  \label{eq:ioudef}
\end{equation}
wobei $|\cdot|$ die Fläche bezeichnet; die zweite Form gilt, weil die
Schnittfläche in $|B_1| + |B_2|$ doppelt gezählt ist. Der Wert liegt zwischen 0
für disjunkte und 1 für deckungsgleiche Boxen.

Wird die \emph{ursprüngliche} Box anschließend als Label der rotierten Instanz
verwendet, so ist das achsparallele Quadrat vollständig in der Hülle enthalten:
Schnittfläche ist $A=d^2$, Vereinigungsfläche $A_{\mathrm{box}} = d^2
(|\cos\theta|+|\sin\theta|)^2$. Mit
$(|\cos\theta|+|\sin\theta|)^{2} = 1 + |\sin 2\theta|$ folgt daraus nach
Gleichung~\eqref{eq:ioudef} die für alle $\theta$ exakte Darstellung
\begin{equation}
  \IoU(\theta) = \frac{1}{1+|\sin 2\theta|},
  \label{eq:iou}
\end{equation}
also $\IoU(2^{\circ})\approx 0{,}935$, $\IoU(10^{\circ})\approx 0{,}745$ und
$\IoU(45^{\circ}) = 0{,}5$.


Die Rechnung stellt einen idealisierten geometrischen Fall dar. Bei kreisförmigen
Rohrstirnflächen ist die tatsächliche Objektgeometrie rotationssymmetrisch und
verändert sich durch eine reine Rotation nicht. Entscheidend ist vielmehr, ob die
verwendete Labeltransformation die Box aus der tatsächlich transformierten
Objektgeometrie neu bestimmt oder lediglich die ursprüngliche achsparallele Box
geometrisch transformiert und anschließend wieder achsparallel umschließt.

Für die vorliegende synthetische Datengenerierung kann diese Unterscheidung
genutzt werden. Werden einzelne Rohrprofile vor der Komposition rotiert, kann
ihre Bounding Box anschließend direkt aus der bekannten Objektmaske bestimmt
werden. Eine geometrische Ungenauigkeit durch die Transformation einer bereits
vorhandenen Box lässt sich dadurch vermeiden.

Aus diesem Grund werden geometrische Variationen, die im Generator label-exakt
umgesetzt werden können, bevorzugt dort erzeugt und nicht zusätzlich als
bildweite Trainingszeit-Transformation angewendet.

\absatz{Kleine Objekte, Auflösung und Labelzuweisung.}\label{subsec:kleine-objekte}
Die Größe eines Objekts im Eingabebild bestimmt gemeinsam mit der Schrittweite
des Detektors, wie viele räumliche Positionen auf einer Merkmalsebene zur
Repräsentation des Objekts zur Verfügung stehen.

Sei $d$ der Durchmesser eines Objekts im Originalbild, $N$ die Kantenlänge des
Originalbildes und $M$ die Kantenlänge der Netzeingabe. Unter der Annahme einer
isotropen Skalierung beträgt die effektive Objektgröße
\begin{equation}
  d_{\mathrm{eff}} = d\,\frac{M}{N}.
  \label{eq:deff}
\end{equation}

Für eine Merkmalsebene mit Stride $S$ beträgt der Abstand benachbarter
Gitterpositionen $S$. Damit ist das Verhältnis
\begin{equation}
  \frac{d_{\mathrm{eff}}}{S}
  \label{eq:ratio}
\end{equation}
ein unmittelbares Maß dafür, wie fein das Objekt durch diese Ebene räumlich
abgetastet wird.

Eine genauere geometrische Betrachtung kann für ein kreisförmiges Objekt
durchgeführt werden. Beide Größen sind dabei im Koordinatensystem der
Netzeingabe zu lesen -- der Stride $S$ ist dort definiert, die Objektgröße muss
also nach Gleichung~\eqref{eq:deff} umgerechnet sein. Unter der idealisierten
Annahme eines regelmäßigen quadratischen Gitters und einer zufälligen,
gleichverteilten Lage des Objekts relativ zu diesem Gitter beträgt die erwartete
Zahl $Z$ der Gitterzentren innerhalb eines Kreises der effektiven Größe
$d_{\mathrm{eff}}$
\begin{equation}
  \mathbb{E}[Z]
  =
  \frac{\pi (d_{\mathrm{eff}}/2)^{2}}{S^{2}}
  =
  \frac{\pi d_{\mathrm{eff}}^{2}}{4S^{2}}.
  \label{eq:gridcount}
\end{equation}
Der Zähler ist die Kreisfläche, der Nenner die Fläche einer Gitterzelle. Weil
jede Zelle genau ein Gitterzentrum enthält, trifft der Kreis über alle Lagen
gemittelt so viele Zentren, wie Zellflächen in seine Fläche passen; bei einer
bestimmten Lage können es mehr oder weniger sein.

Damit steigt die erwartete Zahl räumlicher Positionen quadratisch mit der
Objektgröße und fällt quadratisch mit der Schrittweite. Eine Halbierung der
Schrittweite von $S=8$ auf $S=4$ erhöht die erwartete Zahl von Gitterpositionen
daher um den Faktor vier.

Dies erklärt die Bedeutung einer zusätzlichen P2-Ebene für die vorliegende
Aufgabe. Die P2-Ebene besitzt mit $S=4$ eine doppelt so hohe räumliche Auflösung
wie P3 mit $S=8$. Nach der Dokumentation der Werkzeugkette ist P2 explizit als
Small-Object-Detection-Head vorgesehen \parencite{ultralyticsdocs}; YOLO26
kombiniert diese architektonische Möglichkeit zusätzlich mit STAL, das die
positive Label-Abdeckung kleiner Objekte verbessern soll
\parencite{jocher2026yolo26}.

Die Geometrie des Gitters zeigt außerdem, dass die Existenz einer positiven
Position nicht allein von der Objektgröße abhängt, sondern auch von der
relativen Lage des Objekts zum Gitter. Ein Objekt kann daher selbst dann
ungünstig repräsentiert sein, wenn seine Ausdehnung in der Größenordnung des
Strides liegt.

Für ein kreisförmiges Objekt der effektiven Größe $d_{\mathrm{eff}}$ kann der
Abstand seines Mittelpunkts zum nächstgelegenen Gitterzentrum maximal
$S/\sqrt{2}$ betragen. Ein Gitterzentrum liegt genau dann innerhalb des Kreises,
wenn dieser Abstand kleiner als $d_{\mathrm{eff}}/2$ ist. Damit kann es in einer
idealisierten Betrachtung bereits dann vorkommen, dass kein Gitterzentrum
innerhalb des Objekts liegt, wenn
\begin{equation}
  d_{\mathrm{eff}} < S\sqrt{2}.
  \label{eq:gridmiss}
\end{equation}

Diese Beziehung stellt keine allgemeine Detektionsgrenze dar. Sie beschreibt
lediglich die geometrische Möglichkeit einer fehlenden Gitterposition unter den
genannten Annahmen. Ob ein reales YOLO-Modell ein solches Objekt tatsächlich
nicht detektieren kann, hängt zusätzlich von Feature-Extraktion, Label
Assignment, Regression und Inferenz ab; die Zuweisung bei sehr kleinen Objekten
ist Gegenstand eigener Arbeiten \parencite{zhang2020atss,xu2022rfla}.

Von der Gitterabtastung zu unterscheiden ist ein zweiter, unabhängiger Effekt
der Objektgröße: die Empfindlichkeit der Bewertungsmetrik selbst. Wird eine
quadratische Box der Kantenlänge $d$ um $\delta$ Pixel entlang einer Achse
verschoben, so beträgt die Schnittfläche $(d-\delta)\,d$ und die Vereinigungsfläche
$2d^{2}-(d-\delta)\,d$. Für die Überlappung folgt nach
Gleichung~\eqref{eq:ioudef}
\begin{equation}
  \IoU
  =
  \frac{(d-\delta)\,d}{2d^{2}-(d-\delta)\,d}
  =
  \frac{1-\delta/d}{1+\delta/d}.
  \label{eq:ioushift}
\end{equation}

Maßgeblich ist damit nicht der absolute Boxfehler, sondern das Verhältnis
$\delta/d$. Für $\delta = 2$\px{} -- die Größenordnung des Versatzes, der in
Abschnitt~\ref{sec:boxversatz} zwischen Generatorlabel und Handannotation
gemessen wird -- liefert Gleichung~\eqref{eq:ioushift} bei einem $20$\px{}
großen Objekt $\IoU = 0{,}82$, bei einem $200$\px{} großen dagegen $0{,}98$:
Derselbe Fehler in Pixeln bewirkt beim kleinen Objekt eine rund neunfach
größere Einbuße. Gleichung~\eqref{eq:ioushift} beschreibt eine Verschiebung;
der gemessene Versatz ist dagegen ein Größen-, kein Positionsfehler. Für eine
zentrierte Box, deren Kantenlänge um $\delta$ kleiner ist, gilt
$\IoU = (1-\delta/d)^{2}$, beim $20$\px{} großen Objekt also $0{,}81$. In
erster Ordnung stimmen beide Ausdrücke überein, $\IoU \approx 1 - 2\delta/d$,
sodass die Aussage für beide Fehlerarten gilt. Da die üblichen
Bewertungsschwellen bei $\IoU \geq 0{,}5$ ansetzen und mAP@0,5:0,95 über
Schwellen bis $0{,}95$ mittelt, schlagen wenige Pixel Lokalisierungsfehler bei
kleinen Objekten unmittelbar auf die berichtete Güte durch.

Für die Erzeugung von Trainingsdaten ist das der schärfere der beiden Befunde.
Die Gitterabtastung betrifft die Frage, ob ein Objekt überhaupt eine positive
Position erhält; die Metrikempfindlichkeit betrifft jede einzelne Box, die
erzeugt wird. Bei Objektgrößen im Bereich weniger Dutzend Pixel ist eine
Labelungenauigkeit von ein bis zwei Pixeln deshalb nicht vernachlässigbar,
sondern eine relevante Fehlerquelle; darauf stützt sich die Anforderung an die
Labelgenauigkeit (Abschnitt~\ref{sec:anforderungen}).

Für die vorliegende Untersuchung folgt daraus eine praktische Konsequenz: Die
Trainingsdaten dürfen nicht systematisch Objektgrößen erzeugen, die im Verhältnis
zur gewählten Eingabeauflösung und zum feinsten Detektions-Stride nur
unzureichend repräsentiert werden.

Die Größenverteilung der Trainingsobjekte sollte daher die für die Inferenz
relevante Größenverteilung hinreichend abdecken. Dies betrifft sowohl die
Generierung der Objekte als auch die Wahl der Bildauflösung und der
Skalierungsparameter.

\absatz{Mehrbildverfahren und Mosaic.}\label{subsec:mosaic} Verfahren wie Mosaic kombinieren mehrere Trainingsbilder zu einer gemeinsamen
Eingabe. Dadurch können unterschiedliche Szenen gleichzeitig in einem
Trainingsbeispiel repräsentiert werden und die Varianz der Bildkomposition
erhöht werden.

Für die vorliegende Anwendung ist Mosaic jedoch ambivalent. Einerseits kann das
Verfahren unterschiedliche Hintergründe, Objektpositionen und Objektanordnungen
kombinieren. Andererseits verändert die Zusammensetzung mehrerer Bilder die
effektive Objektgröße. Ein Objekt, das vor der Mosaic-Komposition eine bestimmte
Pixelgröße besitzt, kann durch die anschließende Skalierung deutlich kleiner
werden.

Damit besteht insbesondere bei kleinen Zielobjekten ein Zielkonflikt:
\[
  \text{mehr Szenenvielfalt}
  \quad\longleftrightarrow\quad
  \text{geringere effektive Objektgröße}.
\]

Diese Wechselwirkung ist für die vorliegende Untersuchung relevant, da bereits
die ursprüngliche Aufgabe durch sehr kleine Rohrprofile gekennzeichnet ist.
Mosaic wird daher nicht isoliert als grundsätzlich vorteilhaft oder nachteilig
betrachtet, sondern im Zusammenhang mit der resultierenden
Objektgrößenverteilung.


% ----------------------------------------------------------------------
\section{Rein synthetische Trainingsdaten und Sim-to-Real}
\label{sec:sim2real}

Wenn nahezu der gesamte Trainingsdatensatz synthetisch erzeugt wird, verschiebt
sich die zentrale Fragestellung von der klassischen Datenaugmentation zur
Modellierung einer Trainingsdatenverteilung.

Eine Domäne wird im Folgenden als Paar aus einer Verteilung über dem Eingaberaum
und einer Beschriftungsfunktion verstanden. Sei entsprechend
$\langle \Dom{S}^{X}, f_S\rangle$ die synthetische Quelldomäne und
$\langle \Dom{T}^{X}, f_T\rangle$ die reale Zieldomäne über demselben
Eingaberaum und derselben Aufgabe. Ein Modell $h\in\Hcls$ minimiert im Training
den Fehler auf $\Dom{S}$, soll aber im späteren Einsatz auf $\Dom{T}$ gute
Vorhersagen erzeugen.

Zu betonen ist ferner, dass ein Domänenwechsel eine andere Achse betrifft als
die übliche Aufteilung in Trainings- und Validierungsmenge. Ein Split trennt
zwei Stichproben derselben Verteilung; ein Domänenwechsel bedeutet, dass der
datengenerierende Prozess selbst ein anderer ist.

\absatz{Die Domänenadaptationsschranke als Ordnungsrahmen.}
Für die binäre Klassifikation unter dem $0/1$-Verlust beschränkt die
Domänenadaptationsschranke von \textcite{bendavid2010} den Zielfehler
gleichmäßig für alle $h\in\Hcls$:
\begin{equation}
  \epsilon_T(h)
  \;\leq\;
  \epsilon_S(h)
  \;+\;
  \tfrac{1}{2}\,
  d_{\HdH}\bigl(\Dom{S}^{X},\,\Dom{T}^{X}\bigr)
  \;+\;
  \lambda ,
  \label{eq:bendavid}
\end{equation}
mit dem Quellfehler $\epsilon_S$, dem Zielfehler $\epsilon_T$, der
$\HdH$-Divergenz der beiden Eingabe-Randverteilungen
\[
  d_{\HdH}\bigl(\Dom{S}^{X},\,\Dom{T}^{X}\bigr)
  =
  2 \sup_{h,h'\in\Hcls}
  \Bigl|
    \Pr_{x\sim\Dom{S}^{X}}\bigl[h(x)\neq h'(x)\bigr]
    -
    \Pr_{x\sim\Dom{T}^{X}}\bigl[h(x)\neq h'(x)\bigr]
  \Bigr|
\]
und dem idealen gemeinsamen Fehler
$\lambda = \min_{h'\in\Hcls}\bigl[\epsilon_S(h') + \epsilon_T(h')\bigr]$, den die
beste gemeinsame Hypothese $h^{*}$ erreicht.
Die Schranke macht drei Einflussgrößen sichtbar: wie gut das Modell die synthetische
Domäne lernt, wie stark sich die Eingabeverteilungen beider Domänen
unterscheiden, und ob in $\Hcls$ überhaupt eine Hypothese existiert, die beide
gleichzeitig gut beschreibt.

Zwei Eigenschaften sind für das Folgende wesentlich. Der Divergenzterm ist
ausschließlich über den \emph{Randverteilungen der Eingaben} definiert -- und
bei synthetischen Trainingsdaten ist $\Dom{S}^{X}$ nicht vorgegeben, sondern
durch den Generator \emph{konstruiert}; $p(t)$ ist genau der Freiheitsgrad, über
den dieser Term beeinflusst wird. Merkmale, die der Generator in die Bilder
einführt und die in der Realität fehlen (Tabelle~\ref{tab:signaturen}),
verändern $\Dom{S}^{X}$ und gehören deshalb zum Divergenzterm; ihre
Gegenmaßnahmen sind Teil von $p(t)$. $\lambda$ dagegen hängt von der
Hypothesenklasse ab, ist ohne beschriftete Zieldaten nicht schätzbar, und
Unterschiede zwischen den Beschriftungsfunktionen $f_S \neq f_T$ wirken sich
\emph{ausschließlich} dort aus, in dieser Arbeit etwa die abweichende
Boxkonvention von Generator und Handannotation (Abschnitt~\ref{sec:boxversatz}).

\absatz{Übertragung auf die Detektion.} Gleichung~\eqref{eq:bendavid} setzt
binäre Klassifikation unter dem $0/1$-Verlust voraus. Ihre Struktur überträgt
sich auf Verluste, die symmetrisch sind und die Dreiecksungleichung erfüllen:
An die Stelle von $\tfrac{1}{2}\,d_{\HdH}$ tritt dann die
Discrepancy-Distanz, die für den $0/1$-Verlust mit ihr zusammenfällt, und
$\lambda$ wird mit demselben Verlust gebildet
\parencite{mansour2009discrepancy}. Die Größe $1-\IoU$ erfüllt diese
Bedingungen: Sie ist als Jaccard-Distanz auf Boxen positiver Fläche eine Metrik
\parencite{kosub2019jaccard} und liegt in $[0,1]$. Für ein
Lokalisierungsproblem mit genau einem Objekt je Bild gilt damit eine Schranke
derselben Form, deren Terme mit dem Verlust $1-\IoU$ gebildet sind. Für die
Bewertung dieser Arbeit trägt die Übertragung dennoch nicht:
mAP ist kein Erwartungswert eines Verlusts pro Bild, sondern wird über die nach
Konfidenz geordneten Detektionen des gesamten Datensatzes gebildet, und auch die
Zuordnung einer variablen Zahl von Vorhersagen zu Ground-Truth-Boxen liegt
außerhalb des Rahmens. Gleichung~\eqref{eq:bendavid} wird deshalb \emph{nicht}
als formale Garantie in Anspruch genommen, sondern als begrifflicher
Ordnungsrahmen für die folgenden Entwurfsentscheidungen.

Eine Folgerung trägt dieser Rahmen dennoch: $\lambda$ ist keine Eigenschaft der
Domänen allein, sondern hängt auch von der Bewertung ab. Eine Box gelte als
richtig, sobald ihre IoU zur Referenz mindestens $t_{\IoU}$ beträgt.
Unterscheiden sich $f_S$ und $f_T$ um den in Abschnitt~\ref{sec:boxversatz}
gemessenen Größenversatz von rund zwei Pixeln, so erreicht ein Modell, das $f_S$
exakt reproduziert, an einem $20$\px{} großen Objekt gegenüber $f_T$ nur
$\IoU \approx 0{,}81$ (Abschnitt~\ref{subsec:kleine-objekte}). Für $\lambda$
zählt aber die beste gemeinsame Hypothese $h^{*}$, und sie kann eine Box
mittlerer Größe vorhersagen, rund $19$ statt $18$ oder $20$\px, die zu beiden
Beschriftungen $\IoU \approx 0{,}90$ erreicht. Derselbe Versatz erhöht
$\lambda$ also nur unter strengen Schwellen, hier erst oberhalb von
$t_{\IoU} \approx 0{,}9$: Bei $t_{\IoU} = 0{,}5$ ist er folgenlos, bei
$t_{\IoU} = 0{,}95$ dagegen nicht. Diese Abhängigkeit gilt formal, denn
$\lambda$ ist für jede Bewertung definiert. Nicht gesichert ist unter einer
Schwellenbewertung dagegen die Schranke selbst, weil die Dreiecksungleichung
nicht erfüllt ist. Im Beispiel wird sie für Schwellen zwischen etwa $0{,}81$ und
$0{,}90$ sogar verletzt: Bei gleicher Eingabeverteilung und einem Modell $h$,
das $f_S$ reproduziert, sind $\epsilon_S(h)$, der Divergenzterm und $\lambda$
null, $\epsilon_T(h)$ aber nicht.

\absatz{Sim-to-Real-Abstand.} Ausgewertet wird stattdessen die auf
beschrifteten Realdaten gemessene Differenz
$\Delta_{\mathrm{SR}} = \epsilon_T(h) - \epsilon_S(h)$. Beschriftete reale Daten
sind dafür ausdrücklich nur in der Evaluation erforderlich, nicht im Training.
In Kapitel~\ref{kap:auswertung} wird $\epsilon_S$ auf einem eigenen
synthetischen Split geschätzt, der weder im Training noch bei der Modellauswahl
verwendet wird (\code{sval}, Abschnitt~\ref{sec:simtoreal}). Beide Fehler werden
dort als $\epsilon = 1 - \text{\mAPfifty}$ operationalisiert;
$\Delta_{\mathrm{SR}}$ gibt also an, um wie viel \mAPfifty{} auf einem realen
Split unter dem Wert auf \code{sval} liegt. Aus Gleichung~\eqref{eq:bendavid}
folgt unmittelbar
\begin{equation}
  \Delta_{\mathrm{SR}}
  \;\leq\;
  \tfrac{1}{2}\,
  d_{\HdH}\bigl(\Dom{S}^{X},\,\Dom{T}^{X}\bigr)
  + \lambda .
  \label{eq:gap-bound}
\end{equation}
Für den in \mAPfifty{} gemessenen Abstand ist die Ungleichung aus dem oben
genannten Grund nicht bewiesen, sondern dient nur als Ordnungsrahmen.

Diese Beziehung ist \emph{einseitig}, und die Richtung entscheidet über die
Interpretation. Stimmen beide Domänen in Eingabeverteilung und
Beschriftungsfunktion überein, sind beide Risiken identisch. Ein Abstand, der
über die Stichprobenunsicherheit hinausgeht, belegt daher, dass synthetische und
reale Domäne nicht übereinstimmen -- sei es in der Eingabeverteilung oder in der
Beschriftungsfunktion. Der umgekehrte Schluss ist nicht zulässig: Ein kleiner
Abstand erlaubt keinen Rückschluss auf eine geringe Divergenz, da die Schranke
beliebig locker sein kann und sich die beiden Terme gegenseitig kompensieren
können. $\Delta_{\mathrm{SR}}$ ist folglich kein Maß der Verteilungsdivergenz,
sondern gerichtete Evidenz für Domänenungleichheit. Auch ein sehr niedriger
synthetischer Fehler bei großem Abstand zeigt nur, dass die Summe aus Divergenzterm
und $\lambda$ groß ist, aber nicht, welcher der beiden Terme dafür verantwortlich
ist. In Kapitel~\ref{kap:auswertung} lassen sich beide ihrer Wirkung nach
plausibel zuordnen, ohne dass einer von ihnen gemessen würde: Die abweichende
Boxkonvention ($\lambda$) verschiebt vor allem die strengen IoU-Schwellen
(Abschnitt~\ref{sec:boxversatz}), die fehlende Einbettung (Divergenzterm) kostet
schon bei \mAPfifty{} Rückruf (Abschnitt~\ref{sec:ablationC}). Alle dort
berichteten Werte werden auf endlichen Testmengen berechnet und sind damit
selbst Schätzungen; ihre Stichprobenunsicherheit weist
Abschnitt~\ref{sec:bootstrap} über Konfidenzintervalle aus.

\absatz{Domain Randomization.} Eine verbreitete Strategie gegen einen großen
Sim-to-Real-Abstand besteht darin, die synthetische Domäne absichtlich über die
erwartete reale Variation hinaus zu erweitern. Bei Domain Randomization werden
Parameter des Generators zufällig variiert, sodass das Modell nicht auf eine
einzige, eng definierte Erscheinungsform angewiesen ist \parencite{tobin2017}.
Eine strukturierte Variante beschränkt die Randomisierung dabei auf physikalisch
plausible Parameterbereiche \parencite{prakash2019} -- die Richtung, der auch
Anforderung~A3 folgt.

Randomisiert werden in dieser Arbeit Objektposition und -orientierung,
Objektgröße, Hintergrund, Helligkeit und Kontrast, Rauschen, Unschärfe,
Kompression und Verdeckung; Tabelle~\ref{tab:parameter} im Anhang führt jeden
dieser Parameter mit Bereich und gemessenem Ist-Anteil auf.

Im Rahmen von Gleichung~\eqref{eq:bendavid} ist diese Strategie differenzierter
zu deuten, als es die Vorstellung einer bloßen Verbreiterung nahelegt. Die
$\HdH$-Divergenz misst Unterschiede der Wahrscheinlichkeitsmasse, nicht die
Abdeckung des Trägers. Für jeden Bereich $A = \{x : h(x) \neq h'(x)\}$, auf dem
zwei Hypothesen aus $\Hcls$ uneinig sind, folgt unmittelbar aus der Definition
von $d_{\HdH}$
\begin{equation}
  d_{\HdH}\bigl(\Dom{S}^{X},\,\Dom{T}^{X}\bigr)
  \;\geq\;
  2\bigl(P_T(A) - P_S(A)\bigr),
  \label{eq:massenschranke}
\end{equation}
mit $P_S(A) = \Pr_{x\sim\Dom{S}^{X}}[x \in A]$ und $P_T(A)$ entsprechend.
Lässt sich der Bereich, in dem reale Bilder liegen, durch ein solches $A$
hinreichend genau erfassen, ist $P_T(A) \approx 1$, und jede Masse, die der
Generator außerhalb davon verteilt, fehlt innerhalb davon: Der Divergenzterm
$\tfrac{1}{2}\,d_{\HdH}$ beträgt dann mindestens etwa $1 - P_S(A)$, die Masse
außerhalb. Randomisierung kann diese Untergrenze nur senken, soweit sie Masse
dorthin verschiebt, wo die reale Domäne tatsächlich liegt; über deren Variation
hinaus fortgesetzt, hebt sie die Untergrenze wieder an. Ihr Nutzen ist vor allem
eine Absicherung: Da $\Dom{T}^{X}$ nur unvollständig bekannt ist (worauf auch
Anforderung~A6 beruht), verhindert eine plausible Verbreiterung, dass die
Quelldomäne den realen Bereich verfehlt -- um den Preis, dass der Term nicht
beliebig klein werden kann. Diese Lesart stützt
die strukturierte Variante und das Prinzip aus
Abschnitt~\ref{sec:aug-formalisierung}. Beispiele, die bewusst außerhalb des
realen Bereichs liegen, lassen sich mit der Schranke folglich nicht begründen;
ihr möglicher Nutzen liegt darin, welche Hypothese das Training auswählt, und
darüber macht die Schranke keine Aussage (siehe den Absatz zum Shortcut
Learning).

Hinzu kommt ein zweiter Zielkonflikt: Eine deutlich breitere Quelldomäne
verlangt von $\Hcls$, auf einem größeren Bereich des Eingaberaums gleichzeitig
niedrige Fehler zu erzielen, und kann damit sowohl $\epsilon_S(h)$ als auch
$\lambda$ erhöhen. Die drei Terme sind folglich nicht unabhängig voneinander
optimierbar, was die Notwendigkeit einer experimentellen Bestimmung der
Randomisierungsstärke begründet.

Domain Randomization löst darüber hinaus nicht automatisch das Problem
systematischer Unterschiede zwischen synthetischen und realen Bildern. Wenn eine
bestimmte Eigenschaft in sämtlichen synthetischen Bildern auftritt, bleibt sie
trotz zufälliger Variation anderer Parameter eine mögliche Abkürzung für das
Modell.

\absatz{Shortcut Learning.} Ein Modell kann für eine Aufgabe relevante Vorhersagen über Merkmale treffen,
die zwar im Trainingsdatensatz korreliert, aber nicht kausal oder stabil mit der
Zielklasse verbunden sind. Dieses Verhalten wird als Shortcut Learning
bezeichnet \parencite{geirhos2020}.

Für eine synthetische Datengenerierung entsteht ein besonderes Risiko, da
technische Eigenschaften des Generators systematisch wiederkehren können.
Beispielsweise können harte Objektkanten, identische Schattenmuster oder
charakteristische Farbübergänge in allen synthetischen Bildern auftreten. Ein
Modell kann solche Merkmale nutzen, obwohl sie in realen Bildern nicht vorhanden
sind. In der Terminologie von Gleichung~\eqref{eq:bendavid} ist das ein Beitrag
zum Divergenzterm, nicht zu $\lambda$: Eine Hypothese, die auf die harte
Objektkante reagiert, und eine, die die Stirnfläche selbst erkennt, stimmen auf
synthetischen Bildern überein und widersprechen sich auf realen, und genau
solche Paare machen $d_{\HdH}$ groß. $\lambda$ bleibt klein, weil die zweite in
beiden Domänen richtig liegt; die Schranke sagt nur nicht, welche der beiden das
Training auswählt.

Dabei muss zwischen zufälliger Variation und systematischem Bias unterschieden
werden. Zufällige Variation kann die Vielfalt der beobachteten Beispiele
erhöhen. Ein systematischer Fehler, der in allen synthetischen Bildern vorhanden
ist, wird dagegen durch zusätzliche Ziehungen aus derselben fehlerhaften
Verteilung nicht beseitigt.

Dies ist für die vorliegende Arbeit besonders relevant: Die Qualität der
synthetischen Daten hängt nicht nur davon ab, wie viele verschiedene Bilder
erzeugt werden, sondern auch davon, welche konstanten visuellen Signaturen der
Generator selbst einführt.

\absatz{Synthetische Signaturen.} Für die verwendete Pipeline werden sechs potenzielle Quellen systematischer
Unterschiede zwischen synthetischen und realen Bildern betrachtet. Dass sich
eingefügte Bildbereiche anhand statistischer und lokaler Spuren überhaupt von
ihrer Umgebung unterscheiden lassen, ist Gegenstand der Bildforensik
\parencite{farid2009}; die dort untersuchten Merkmale sind zugleich diejenigen,
die ein Faltungsnetz als Abkürzung verwenden kann.

\begin{table}[htbp]
  \centering
  \caption{Betrachtete synthetische Signaturen und zugehörige Gegenmaßnahmen.}
  \label{tab:signaturen}
  \small
  \begin{tabularx}{\linewidth}{@{}l X X@{}}
    \toprule
    \textbf{Signatur} & \textbf{Ursache} & \textbf{Mögliche Gegenmaßnahme} \\
    \midrule
    S1 & sichtbare Übergänge an eingefügten Objekten
       & weiche Masken, Farbangleichung \\
    S2 & künstliche Helligkeits- oder Farbverhältnisse
         \parencite{reinhard2001,tsai2017}
       & Variation und Matching photometrischer Parameter \\
    S3 & unrealistische Schärfe oder Auflösung; Skalierung hinterlässt
         messbare Resampling-Spuren \parencite{popescu2005,zhangr2019}
       & Blur- und Downsampling-Variation \\
    S4 & fehlendes bzw.\ unrealistisches Sensorrauschen; Rauschmuster sind
         kameraspezifisch \parencite{lukas2006}
       & Rauschmodell \\
    S5 & fehlende Kompressionsartefakte; eingefügte Pixel tragen eine andere
         DCT-Blockstatistik \parencite{wallace1992jpeg}
       & JPEG-Kompression \\
    S6 & fehlende Interaktion mit der Umgebung
       & physikalisch plausible Okkluder und Schatten \\
    \bottomrule
  \end{tabularx}
\end{table}

Tabelle~\ref{tab:signaturen} stellt keine vollständige Beschreibung aller
möglichen Domain-Gap-Ursachen dar. Sie definiert vielmehr die sechs Signaturen,
die für die in dieser Arbeit entwickelte Pipeline als besonders relevant
betrachtet und experimentell untersucht werden.

\absatz{Verdeckung als Domäneneigenschaft.} Bei den untersuchten Rohrprofilen ist Verdeckung ein normaler Bestandteil der
Zielszene. Ein Rohrende kann teilweise durch davorliegende Profile, das Gitter
oder andere Bildelemente verdeckt werden.

Daraus ergibt sich ein Unterschied zu klassischen Verfahren wie Cutout oder
Random Erasing. Dort wird Verdeckung primär als künstliche Störung eingesetzt. Im
vorliegenden Fall soll dagegen die reale Ursache der Verdeckung nachgebildet
werden.

Die Qualität eines Okkluders besitzt deshalb zwei Dimensionen:
\begin{enumerate}[label=\arabic*.,itemsep=2pt]
  \item \textbf{geometrische Wirkung:} Wie viel des Zielobjekts wird verdeckt?
  \item \textbf{visuelle Plausibilität:} Wie ähnlich ist der Okkluder einem
        tatsächlich vor dem Objekt liegenden Profilelement?
\end{enumerate}

Insbesondere die zweite Dimension kann relevant werden, wenn das Modell nicht
nur das sichtbare Rohrende, sondern auch die charakteristischen visuellen
Übergänge zwischen Objekt und Okkluder als Erkennungsmerkmal verwendet.

Ein künstlich erzeugter Okkluder kann dadurch selbst zu einer unerwünschten
Trainingssignatur werden. Werden beispielsweise alle Verdeckungen mit identischen
Kanten, Farben oder Schatten erzeugt, könnte das Modell lernen, diese künstlichen
Muster mit der Zielklasse zu verbinden.

Die Hypothese dieser Arbeit lautet daher, dass nicht allein die Menge der
Verdeckung, sondern auch deren visuelle Realistik die Übertragbarkeit auf reale
Bilder beeinflusst. Diese Hypothese wird in der späteren Ablation experimentell
überprüft.

% ----------------------------------------------------------------------
\section{Anforderungen an die Datengenerierungspipeline}
\label{sec:anforderungen}

Aus den vorangegangenen Überlegungen ergeben sich sechs Anforderungen an die in
dieser Arbeit entwickelte Pipeline.

\begin{description}[leftmargin=0pt,labelsep=0.5em,itemsep=3pt,
                    font=\normalfont\bfseries]
  \item[A1 -- Labelgenauigkeit.] Geometrische Transformationen müssen mit
  möglichst geringer Abweichung zwischen Objekt und Bounding Box umgesetzt
  werden; insbesondere darf keine zusätzliche Labelungenauigkeit entstehen, wo
  die Objektgeometrie im Generator bekannt ist.

  \item[A2 -- Kalibrierte Objektgrößen.] Die erzeugte Größenverteilung muss die
  für die reale Inferenz relevante abdecken, ohne durch Skalierung systematisch
  Objekte zu erzeugen, deren effektive Größe im Verhältnis zu Eingabeauflösung
  und feinstem Detektions-Stride zu klein ist.

  \item[A3 -- Physikalisch plausible Größenabhängigkeiten.] Parameter, die in
  der realen Szene voneinander abhängen, dürfen nicht unabhängig randomisiert
  werden -- Objektgröße, perspektivische Darstellung und die Größe gemeinsam
  auftretender Instanzen sind konsistent zu behandeln.

  \item[A4 -- Kontrolle synthetischer Signaturen.] Für jede der Signaturen
  S1--S6 muss eine Gegenmaßnahme vorgesehen und ihre Wirksamkeit, soweit der
  Validierungsdatensatz das zulässt, experimentell geprüft werden.

  \item[A5 -- Kontrollierbare Variation.] Die Verteilung $p(t)$ der
  Generatorparameter muss explizit dokumentiert und reproduzierbar sein und eine
  begründete Approximation der erwarteten Variation der Zieldomäne bilden.

  \item[A6 -- Trennung von Annahme und Messung.] Eigenschaften der realen
  Domäne, die nicht gemessen werden können, sind als Annahmen zu kennzeichnen;
  Parameter, die allein auf einer Plausibilitätsüberlegung beruhen, dürfen nicht
  als empirisch validiert dargestellt werden.
\end{description}

Die theoretische Analyse bestimmt damit nicht die optimale Augmentierung a
priori. Sie definiert die Eigenschaften, die eine geeignete Pipeline erfüllen
sollte, und trennt so drei Ebenen der Untersuchung: theoretische Anforderungen,
synthetische Datengenerierung, experimentelle Validierung. Ob die abgeleiteten
Entscheidungen tatsächlich zu einer besseren Übertragung auf reale Bilder
führen, wird in Kapitel~\ref{kap:auswertung} empirisch untersucht.


%\end{document}

% ======================================================================
%  Optional, falls spaeter Alignment-Verfahren (z.B. DANN) diskutiert
%  werden; dieser Eintrag steht noch nicht in literatur.bib.
% ======================================================================
%
% @inproceedings{zhao2019invariant,
%   author    = {Zhao, Han and Combes, Remi Tachet des and Zhang, Kun and
%                Gordon, Geoffrey J.},
%   title     = {On learning invariant representations for domain adaptation},
%   booktitle = {Proceedings of the 36th International Conference on
%                Machine Learning (ICML)},
%   year      = {2019}
% }
